\scanpage{11}
On voit ainsi que les propriétés des actions de $\Gamma,\deleted{\unsure{\Gamma'}}$ sur les ens de sommets, arrêtes et faces de l'icosaèdre (\deleted{ensembles} munis de leurs relations d'incidence mutuelles) sont conséquences de la seule structure combinatoire de ces ensembles (soit \deleted{Il convient} sommets-arrêtes, soit faces-arrêtes, soit \unsure{mieux|même} le triple système sommets-arrêtes-faces \ldots). Il convient cependant de caractériser le sous-groupe $\Gamma'$ de $\Gamma$ en termes combinatoires, On peut dire p. ex. que c'est le groupe qui respecte \inserted{les} ordres circulaires qu'on a sur les \inserted{ens. des points périphériques des} calottes \deleted{sphériques} de \deleted{\unsure{l'icosa}} l'icosaèdre (en termes de la structure combinatoire, on peut définir \uline{deux} ``syst. cohérents'' d'ordres circulaires sur les calottes \deleted{\unsure{sphér}} $\inserted{\unsure{\pm}}$, correspondants aux deux orientations possibles de l'espace ambiant).
% archive sheet number: 10
